\scanpage{165}
\begin{center}- 46 -\end{center}\hfill IX

projectif implique que $\alpha_s(F_K)$ se comporte additivement pour les suites exactes.

Nous allons appliquer cette définition au cas du faisceau ${}_\ell A_K$, où $A_K$ est notre schéma abélien sur $K$. On a donc, dans ce cas, un entier
\[
 (4.5.4)\quad\alpha_s(A_K;\ell)=\operatorname{long}\operatorname{Hom}_{\mathbb Z_\ell(\Gamma)}(\mathrm{Sw},U_\ell\otimes_{\mathbb Z_\ell}\mathbb F_\ell),\quad U_\ell=T_\ell(A(\overline K)),
\]
$\Gamma$ étant choisi comme il a été dit (choix dépendant à priori de $\ell$).
L'entier (4.5.4) s'appelle l'\uline{exposant conducteur immodéré} (ou ``\uline{sauvage}'') (*) du schéma abélien $A_K$ par rapport à l'anneau de valuation discrète $V\subset K$ ; par passage au hensélisé, il est défini bien sûr sans supposer $V$ hensélien, et il se ramène aussitôt au cas où $S$ est \uline{strictement} local. Nous nous proposons de montrer le

\uline{Corollaire} 4.6. \uline{L'exposant conducteur immodéré} (4.5.4) \uline{ne dépend pas du choix du nombre premier} $\ell\ne p$.
\figure{\begin{tikzpicture}\draw (0,.12)--(.10,0);\draw (.03,.12)--(.13,0);\end{tikzpicture}}

Nous pouvons supposer $S$ strictement local. Nous allons donner le calcul de (4.5.4), en remplaçant $\mathrm{Sw}$ par n'importe quel autre $\mathbb Z_\ell(\Gamma)$-module projectif $L$. Notons que la filtration (4.1.1) de $U_\ell$ se fait par des sous-$\mathbb Z_\ell$-modules qui sont manifestement des facteurs directs, donc

\figure{\begin{tikzpicture}\draw (0,0)--(3,0);\end{tikzpicture}}
(*) Il y a lieu de définir aussi l'\uline{exposant conducteur modéré} de $A_K$ par rapport à $V$ comme l'entier différence des longueurs sur $\mathbb Z/\ell\mathbb Z$ des fibres géométriques générique et spéciale de ${}_\ell A^{\unsure{o|0}}$, qui est indépendant de $\ell$ et égal à $2\lambda+\mu$ avec les notations de (1.1.11) ; \uline{l'exposant conducteur} (total) de $A_K$ est défini comme la somme des deux entiers précédents; c'est le terme local qui intervient dans la formule d'Euler-Poincaré pour ${}_\ell A^{\unsure{o|0}}$ de SGA 5 X 6.2, [20].
\margin{$\unsure{\Lambda|A}$\quad\unsure{214|211}}
\margin{364}
\margin{164}
